\chapter{Notação utilizada}

    \begin{table}[h!]
    \centering
    \begin{tabular}{ >{\raggedleft\arraybackslash}p{2.2cm} |  >{\arraybackslash}p{5cm} >{\raggedleft\arraybackslash}p{2.2cm} |  >{\arraybackslash}p{5cm} }\hline
        \multicolumn{4}{c}{\textbf{Variáveis termodinâmicas}}\\\hline
        \textbf{Notação} & \textbf{Descrição} & \textbf{Notação} & \textbf{Descrição}\\\hline
        $U$ & Energia interna & $T$ & Temperatura\\
        $N$ & Número de partículas & $\mu$ & Potencial químico\\
        $V$ & Volume (tamanho) & $p$ & Pressão\\
        $F$ & Energia livre de Helmholtz & $G$ & Energia livre de Gibbs\\
        $H$ & Entalpia & $\mathcal{A}$ & Grande potencial\\
        $S$ & Entropia & $\pps$& Entropia por partícula\\
        \hline
        \multicolumn{4}{c}{\textbf{Variáveis estatísticas}}\\\hline
        $\Omega(U,N,V)$ & \multicolumn{3}{l}{Multiplicidade do macroestado $(U,N,V)$}\\
        $\pzcal{Z}_C(T,N,V)$ & \multicolumn{3}{l}{Função de partição canônica do estado $(T,N,V)$}\\
        $z_1(T,V)$ & \multicolumn{3}{l}{Função de partição canônica de partícula única no estado $(T,V)$}\\
        $\mathcal{Z}_G(T,\mu,V)$ & \multicolumn{3}{l}{Função de partição grande canônica do estado $(T,\mu,V)$}\\
        $E(j)$ & \multicolumn{3}{l}{Energia associada ao estado $j$ do ensemble}\\
        $N(j)$ & \multicolumn{3}{l}{Número de partículas associada ao estado $j$ do ensamble}\\
    \end{tabular}
    \caption{Notação padrão utilizada}
    \label{tab:notacao_padrao_ensembles_estatisticos}
    \end{table}